% !PRINTED 50
\weberpage{50}
%<<CONT

\[
F(x)=x+3{x(x-1)\over1\cdot2}+2{x(x-1)(x-2)\over1\cdot2\cdot3}
={x(x+1)(2x+1)\over6}.
\]
Für \(f(x)=x^3\) ergiebt dieselbe Rechnung
\[
F(x)=\left({x(x+1)\over2}\right)^2.
\]
Die Summe der \(m\) ersten Cuben ist also gleich dem Quadrat der \(m\)ten Trigonalzahl.

\begin{center}
{\large \S.\ 12.}

\vspace{0.5em}
{\large\gesperrt{Der polynomische Lehrsatz.}}
\end{center}

Im \S~8 ist für den binomischen Lehrsatz die Form abgeleitet:
\begin{equation}
 (x+y)^n=\Pi(n)\sum^{\alpha,\beta}{x^\alpha y^\beta\over \Pi(\alpha)\Pi(\beta)},
\tag{1}
\end{equation}
in der sich die Summe auf alle Combinationen zweier Zahlen $\alpha,\beta$ erstreckt, deren keine negativ ist und die der Bedingung
\begin{equation}
 \alpha+\beta=n
\tag{2}
\end{equation}
genügen.

Diese Form gestattet, zunächst durch Induction, eine Verallgemeinerung auf die $n$te Potenz eines Polynoms:
\begin{equation}
 (x+y+z+\cdots)^n=\Pi(n)\sum^{\alpha,\beta,\gamma,\ldots}
 {x^\alpha y^\beta z^\gamma\cdots
\over \Pi(\alpha)\Pi(\beta)\Pi(\gamma)\cdots},
\tag{3}
\end{equation}
mit der Bestimmung, dass $\alpha,\beta,\gamma,\ldots$ alle positiven oder verschwindenden ganzzahligen Werthe durchlaufen, die der Bedingung
\begin{equation}
 \alpha+\beta+\gamma+\cdots=n
\tag{4}
\end{equation}
genügen. Um aber die Richtigkeit dieser Formel allgemein zu beweisen, nehmen wir an, sie sei bewiesen, wenn das Polynom ein Glied weniger enthält, wie sie es in der That ist, wenn das Polynom nur zwei Glieder enthält.

Wir setzen dann
\begin{equation}
 u=y+z+\cdots
\tag{5}
\end{equation}
und wenden auf $(x+u)$ die Formel (1) an, aus der sich ergiebt:
%CONT>>
